\documentclass[12pt,oneside,a4paper]{book} % This style for A4 format.

\input{./AuxFiles/Packages}

%%%%%%%%%%%%%%%%%%%%%%%%%%%%%%%%%%%%%%%%%%%%%%%%%%%%%%%%%%%%%%%%
%% Document settings
%%%%%%%%%%%%%%%%%%%%%%%%%%%%%%%%%%%%%%%%%%%%%%%%%%%%%%%%%%%%%%%%

\title{Comparative Analysis of Evapotranspiration and Crop Yield in Hemavati Left Bank Canal Command Area using QGIS and Remote Sensing}

\subCode{22CVMP}

\Department[CV]{Civil Engineering}

\academicYear{2022-2026}

%%%%%%%%%%%%%%%%%%%%%%%%%%%%%%%%%%%%%%%%%%%%%%%%%%%%%%%%%%%%%%%%
%% Report Type
%%%%%%%%%%%%%%%%%%%%%%%%%%%%%%%%%%%%%%%%%%%%%%%%%%%%%%%%%%%%%%%%

% UG Major Project
% Keep this commented for Major Project

%\MinorProject

%%%%%%%%%%%%%%%%%%%%%%%%%%%%%%%%%%%%%%%%%%%%%%%%%%%%%%%%%%%%%%%%
%% Student Details
%%%%%%%%%%%%%%%%%%%%%%%%%%%%%%%%%%%%%%%%%%%%%%%%%%%%%%%%%%%%%%%%

\stuNameA{Brunda K S}
\stuUSNA{1RV22CV017}

\stuNameB{Deepti K M}
\stuUSNB{1RV22CV024}

\stuNameC{Suprith S N}
\stuUSNC{1RV22CV110}

\stuNameD{Srusthi Ashok Kosthi}
\stuUSND{1RV23CV415}

%%%%%%%%%%%%%%%%%%%%%%%%%%%%%%%%%%%%%%%%%%%%%%%%%%%%%%%%%%%%%%%%
%% Guide Related Commands
%%%%%%%%%%%%%%%%%%%%%%%%%%%%%%%%%%%%%%%%%%%%%%%%%%%%%%%%%%%%%%%%

\guideNameA{Mr. Gowtham Prasad M E}
\guideDesignationA{Assistant Professor}
\guideDeptA[CV]{Civil Engineering}
\guideOrgA{RV College of Engineering}

%%%%%%%%%%%%%%%%%%%%%%%%%%%%%%%%%%%%%%%%%%%%%%%%%%%%%%%%%%%%%%%%
%% External Guide (Optional)
%%%%%%%%%%%%%%%%%%%%%%%%%%%%%%%%%%%%%%%%%%%%%%%%%%%%%%%%%%%%%%%%

%\guideNameB{}
%\guideDesignationB{}
%\guideDeptB{}
%\guideOrgB{}

%%%%%%%%%%%%%%%%%%%%%%%%%%%%%%%%%%%%%%%%%%%%%%%%%%%%%%%%%%%%%%%%
%% Panel Members
%%%%%%%%%%%%%%%%%%%%%%%%%%%%%%%%%%%%%%%%%%%%%%%%%%%%%%%%%%%%%%%%

\panelMemberNameA{Dr. XYZ}
\panelMemberDesigA{Professor}
\panelMemberDeptA[CV]{Civil Engineering}

\panelMemberNameB{Prof. ABC}
\panelMemberDesigB{Assistant Professor}
\panelMemberDeptB[CV]{Civil Engineering}

%%%%%%%%%%%%%%%%%%%%%%%%%%%%%%%%%%%%%%%%%%%%%%%%%%%%%%%%%%%%%%%%
%% Project Coordinators
%%%%%%%%%%%%%%%%%%%%%%%%%%%%%%%%%%%%%%%%%%%%%%%%%%%%%%%%%%%%%%%%

\projectCoOrdNameA{Dr. XYZ}
\projectCoOrdDesigA{Professor}
\projectCoOrdDeptA[CV]{Civil Engineering}

\projectCoOrdNameB{Dr. ABC}
\projectCoOrdDesigB{Professor}
\projectCoOrdDeptB[CV]{Civil Engineering}

%%%%%%%%%%%%%%%%%%%%%%%%%%%%%%%%%%%%%%%%%%%%%%%%%%%%%%%%%%%%%%%%
%% Esteemed Members
%%%%%%%%%%%%%%%%%%%%%%%%%%%%%%%%%%%%%%%%%%%%%%%%%%%%%%%%%%%%%%%%

\HOD{Dr. HOD NAME}

\Principal{Dr. K. N. Subramanya}

%%%%%%%%%%%%%%%%%%%%%%%%%%%%%%%%%%%%%%%%%%%%%%%%%%%%%%%%%%%%%%%%
%% QR URL
%%%%%%%%%%%%%%%%%%%%%%%%%%%%%%%%%%%%%%%%%%%%%%%%%%%%%%%%%%%%%%%%

\QRurl{}

%%%%%%%%%%%%%%%%%%%%%%%%%%%%%%%%%%%%%%%%%%%%%%%%%%%%%%%%%%%%%%%%
%% Draft Report
%%%%%%%%%%%%%%%%%%%%%%%%%%%%%%%%%%%%%%%%%%%%%%%%%%%%%%%%%%%%%%%%

%\DraftCopy

%%%%%%%%%%%%%%%%%%%%%%%%%%%%%%%%%%%%%%%%%%%%%%%%%%%%%%%%%%%%%%%%
%% Glossaries & Acronyms
%%%%%%%%%%%%%%%%%%%%%%%%%%%%%%%%%%%%%%%%%%%%%%%%%%%%%%%%%%%%%%%%

\newglossary[sym]{symbolList}{sym1}{sym2}{List of Symbols}

\makeglossaries

\loadglsentries{./AuxFiles/Glossaries}

\renewcommand{\glspostdescription}{}

%%%%%%%%%%%%%%%%%%%%%%%%%%%%%%%%%%%%%%%%%%%%%%%%%%%%%%%%%%%%%%%%
%% Bibliography
%%%%%%%%%%%%%%%%%%%%%%%%%%%%%%%%%%%%%%%%%%%%%%%%%%%%%%%%%%%%%%%%

\usepackage[backend=bibtex,style=ieee]{biblatex}

\addbibresource{./AuxFiles/ProjectBib.bib}

%%%%%%%%%%%%%%%%%%%%%%%%%%%%%%%%%%%%%%%%%%%%%%%%%%%%%%%%%%%%%%%%
%% WaterMark
%%%%%%%%%%%%%%%%%%%%%%%%%%%%%%%%%%%%%%%%%%%%%%%%%%%%%%%%%%%%%%%%

\usepackage{background}

\backgroundsetup{
scale=1,
angle=0,
firstpage = false,
opacity=0.5,
contents={
\begin{tikzpicture}[remember picture, overlay]
\node at ([yshift=0pt, xshift=0pt]current page.center)
{\includegraphics[width=0.6\textwidth]{Figures/RVlogoVecW}};
\end{tikzpicture}
}
}

%%%%%%%%%%%%%%%%%%%%%%%%%%%%%%%%%%%%%%%%%%%%%%%%%%%%%%%%%%%%%%%%
%% Begin Document
%%%%%%%%%%%%%%%%%%%%%%%%%%%%%%%%%%%%%%%%%%%%%%%%%%%%%%%%%%%%%%%%

\begin{document}

\NoBgThispage

\maketitle

\newpage

%%%%%%%%%%%%%%%%%%%%%%%%%%%%%%%%%%%%%%%%%%%%%%%%%%%%%%%%%%%%%%%%
%% Certificate
%%%%%%%%%%%%%%%%%%%%%%%%%%%%%%%%%%%%%%%%%%%%%%%%%%%%%%%%%%%%%%%%

\chapter*{Certificate}

This is to certify that the project work entitled:

\begin{center}

\textbf{``Comparative Analysis of Evapotranspiration and Crop Yield in Hemavati Left Bank Canal Command Area using QGIS and Remote Sensing''}

\end{center}

is a bonafide work carried out by:

\begin{itemize}

\item Brunda K S (1RV22CV017)

\item Deepti K M (1RV22CV024)

\item Suprith S N (1RV22CV110)

\item Srusthi Ashok Kosthi (1RV23CV415)

\end{itemize}

in partial fulfillment of the requirements for the award of Bachelor of Engineering in Civil Engineering of Visvesvaraya Technological University, Belagavi during the academic year 2022--2026.

\vspace{2cm}

\begin{tabular}{p{7cm}p{7cm}}

Guide Signature & HOD Signature \\

\\

Mr. Gowtham Prasad M E &
Head of Department \\

Assistant Professor &
Department of Civil Engineering \\

\end{tabular}

%%%%%%%%%%%%%%%%%%%%%%%%%%%%%%%%%%%%%%%%%%%%%%%%%%%%%%%%%%%%%%%%
%% Declaration
%%%%%%%%%%%%%%%%%%%%%%%%%%%%%%%%%%%%%%%%%%%%%%%%%%%%%%%%%%%%%%%%

\newpage

\chapter*{Declaration}

We hereby declare that the project entitled:

\begin{center}

\textbf{``Comparative Analysis of Evapotranspiration and Crop Yield in Hemavati Left Bank Canal Command Area using QGIS and Remote Sensing''}

\end{center}

submitted to RV College of Engineering is a record of original work carried out by us under the guidance of Mr. Gowtham Prasad M E.

\vspace{2cm}

\begin{flushright}

Brunda K S \\

Deepti K M \\

Suprith S N \\

Srusthi Ashok Kosthi

\end{flushright}

%%%%%%%%%%%%%%%%%%%%%%%%%%%%%%%%%%%%%%%%%%%%%%%%%%%%%%%%%%%%%%%%
%% Acknowledgement
%%%%%%%%%%%%%%%%%%%%%%%%%%%%%%%%%%%%%%%%%%%%%%%%%%%%%%%%%%%%%%%%

\newpage

\chapter*{Acknowledgement}

We express our sincere gratitude to RV College of Engineering, Bengaluru, for providing us the opportunity to carry out this major project.

We are deeply thankful to our guide Mr. Gowtham Prasad M E, Assistant Professor, Department of Civil Engineering, for his constant support, encouragement, and valuable suggestions throughout the project.

We also thank the faculty members of the Department of Civil Engineering for their support and guidance.

Finally, we express our gratitude to our family and friends for their encouragement during the completion of this work.

%%%%%%%%%%%%%%%%%%%%%%%%%%%%%%%%%%%%%%%%%%%%%%%%%%%%%%%%%%%%%%%%
%% Abstract
%%%%%%%%%%%%%%%%%%%%%%%%%%%%%%%%%%%%%%%%%%%%%%%%%%%%%%%%%%%%%%%%

\newpage

\pagenumbering{roman}

\chapter*{Abstract}

Efficient irrigation management is essential for sustainable agricultural development and effective water resource utilization. This project focuses on the comparative analysis of evapotranspiration and crop yield in the Hemavati Left Bank Canal command area using Remote Sensing and GIS techniques.

Landsat-8 and Sentinel-2 satellite imagery were utilized for generating vegetation indices such as NDVI and EVI. Actual evapotranspiration maps were generated using the Penman-Monteith equation and crop coefficient values derived from FAO standards.

QGIS software was used for spatial analysis, thematic map generation, and crop water requirement estimation. Paddy and ragi crops were assumed to be dominant crops cultivated within the command area.

The generated NDVI, EVI, ET, and Kc maps were analyzed to study spatial variations in vegetation and irrigation demand. The relationship between evapotranspiration and crop yield was evaluated to identify irrigation efficiency across the study region.

The study demonstrates the effectiveness of GIS and Remote Sensing techniques in agricultural water management and crop monitoring.

%%%%%%%%%%%%%%%%%%%%%%%%%%%%%%%%%%%%%%%%%%%%%%%%%%%%%%%%%%%%%%%%
%% Table of Contents
%%%%%%%%%%%%%%%%%%%%%%%%%%%%%%%%%%%%%%%%%%%%%%%%%%%%%%%%%%%%%%%%

\tableofcontents

\newpage

\listoffigures

\newpage

\listoftables

%%%%%%%%%%%%%%%%%%%%%%%%%%%%%%%%%%%%%%%%%%%%%%%%%%%%%%%%%%%%%%%%
%% Main Matter
%%%%%%%%%%%%%%%%%%%%%%%%%%%%%%%%%%%%%%%%%%%%%%%%%%%%%%%%%%%%%%%%

\mainmatter

%%%%%%%%%%%%%%%%%%%%%%%%%%%%%%%%%%%%%%%%%%%%%%%%%%%%%%%%%%%%%%%%
%% Chapter 1
%%%%%%%%%%%%%%%%%%%%%%%%%%%%%%%%%%%%%%%%%%%%%%%%%%%%%%%%%%%%%%%%

\chapter{Introduction}

\section{General}

Agriculture is one of the major water-consuming sectors in India. Efficient irrigation management is necessary to improve agricultural productivity and sustainable water utilization.

Remote Sensing and GIS techniques provide effective tools for monitoring crop conditions, evapotranspiration, and crop water requirements over large agricultural areas.

\section{Need for Study}

The Hemavati Left Bank Canal command area supports extensive agricultural activities. Proper estimation of evapotranspiration and crop water requirement helps improve irrigation management and crop productivity.

\section{Objectives}

\begin{itemize}

\item To generate NDVI and EVI maps using satellite imagery.

\item To estimate actual evapotranspiration using Remote Sensing techniques.

\item To develop crop coefficient maps using FAO values.

\item To estimate crop water requirements for paddy and ragi crops.

\item To compare evapotranspiration and crop yield.

\end{itemize}

%%%%%%%%%%%%%%%%%%%%%%%%%%%%%%%%%%%%%%%%%%%%%%%%%%%%%%%%%%%%%%%%
%% Chapter 2
%%%%%%%%%%%%%%%%%%%%%%%%%%%%%%%%%%%%%%%%%%%%%%%%%%%%%%%%%%%%%%%%

\chapter{Literature Review}

Several studies have demonstrated the use of Remote Sensing and GIS techniques in irrigation management and agricultural monitoring.

NDVI and EVI are widely used vegetation indices for analyzing crop health and vegetation density. The Penman-Monteith method is extensively used for evapotranspiration estimation.

Remote sensing provides accurate spatial information regarding crop conditions, vegetation distribution, and water demand in agricultural regions.

GIS tools assist in integrating satellite imagery, evapotranspiration data, and crop yield information for irrigation planning and sustainable agricultural management.

%%%%%%%%%%%%%%%%%%%%%%%%%%%%%%%%%%%%%%%%%%%%%%%%%%%%%%%%%%%%%%%%
%% Chapter 3
%%%%%%%%%%%%%%%%%%%%%%%%%%%%%%%%%%%%%%%%%%%%%%%%%%%%%%%%%%%%%%%%

\chapter{Study Area}

The study area selected for this project is the Hemavati Left Bank Canal command area located in Karnataka.

The region supports cultivation of major crops such as paddy and ragi. Irrigation water is supplied through the Hemavati canal system.

Landsat-8 and Sentinel-2 imagery were used for the study to generate NDVI, EVI, ET, and Kc maps.

%%%%%%%%%%%%%%%%%%%%%%%%%%%%%%%%%%%%%%%%%%%%%%%%%%%%%%%%%%%%%%%%
%% Study Area Figure
%%%%%%%%%%%%%%%%%%%%%%%%%%%%%%%%%%%%%%%%%%%%%%%%%%%%%%%%%%%%%%%%

\begin{figure}[H]

\centering

\includegraphics[width=0.95\textwidth]{Figures/Study_Area_Map.png}

\caption{Study Area Map}

\label{fig:studyarea}

\end{figure}

%%%%%%%%%%%%%%%%%%%%%%%%%%%%%%%%%%%%%%%%%%%%%%%%%%%%%%%%%%%%%%%%
%% Chapter 4
%%%%%%%%%%%%%%%%%%%%%%%%%%%%%%%%%%%%%%%%%%%%%%%%%%%%%%%%%%%%%%%%

\chapter{Methodology}

\section{Satellite Data Used}

\begin{itemize}

\item Landsat-8 imagery

\item Sentinel-2 imagery

\end{itemize}

\section{Software Used}

\begin{itemize}

\item QGIS

\item Remote Sensing tools

\end{itemize}

%%%%%%%%%%%%%%%%%%%%%%%%%%%%%%%%%%%%%%%%%%%%%%%%%%%%%%%%%%%%%%%%
%% NDVI
%%%%%%%%%%%%%%%%%%%%%%%%%%%%%%%%%%%%%%%%%%%%%%%%%%%%%%%%%%%%%%%%

\section{NDVI Calculation}

\[
NDVI = \frac{NIR - Red}{NIR + Red}
\]

NDVI was used to analyze vegetation density and crop health in the study area.

\begin{figure}[H]

\centering

\includegraphics[width=0.95\textwidth]{Figures/NDVI_Map.png}

\caption{NDVI Map of Hemavati Left Bank Canal Command Area}

\label{fig:ndvi}

\end{figure}

%%%%%%%%%%%%%%%%%%%%%%%%%%%%%%%%%%%%%%%%%%%%%%%%%%%%%%%%%%%%%%%%
%% EVI
%%%%%%%%%%%%%%%%%%%%%%%%%%%%%%%%%%%%%%%%%%%%%%%%%%%%%%%%%%%%%%%%

\section{EVI Calculation}

\[
EVI = G \times \frac{NIR - Red}{NIR + C_1(Red) - C_2(Blue) + L}
\]

EVI was used to improve vegetation analysis in dense vegetation regions.

\begin{figure}[H]

\centering

\includegraphics[width=0.95\textwidth]{Figures/EVI_Map.png}

\caption{EVI Map of Hemavati Left Bank Canal Command Area}

\label{fig:evi}

\end{figure}

%%%%%%%%%%%%%%%%%%%%%%%%%%%%%%%%%%%%%%%%%%%%%%%%%%%%%%%%%%%%%%%%
%% Evapotranspiration
%%%%%%%%%%%%%%%%%%%%%%%%%%%%%%%%%%%%%%%%%%%%%%%%%%%%%%%%%%%%%%%%

\section{Evapotranspiration Estimation}

The Penman-Monteith equation was used for estimating actual evapotranspiration.

\[
ET_0 =
\frac{
0.408\Delta(R_n-G)+\gamma\frac{900}{T+273}u_2(e_s-e_a)
}{
\Delta+\gamma(1+0.34u_2)
}
\]

\begin{figure}[H]

\centering

\includegraphics[width=0.95\textwidth]{Figures/ET_Map.png}

\caption{Actual Evapotranspiration Map}

\label{fig:et}

\end{figure}

%%%%%%%%%%%%%%%%%%%%%%%%%%%%%%%%%%%%%%%%%%%%%%%%%%%%%%%%%%%%%%%%
%% KC MAP
%%%%%%%%%%%%%%%%%%%%%%%%%%%%%%%%%%%%%%%%%%%%%%%%%%%%%%%%%%%%%%%%

\section{Crop Coefficient Mapping}

Crop coefficient values were assumed using FAO standards for paddy and ragi crops.

\[
ET_c = K_c \times ET_0
\]

\begin{figure}[H]

\centering

\includegraphics[width=0.95\textwidth]{Figures/Kc_Map.png}

\caption{Crop Coefficient (Kc) Map}

\label{fig:kc}

\end{figure}

%%%%%%%%%%%%%%%%%%%%%%%%%%%%%%%%%%%%%%%%%%%%%%%%%%%%%%%%%%%%%%%%
%% Chapter 5
%%%%%%%%%%%%%%%%%%%%%%%%%%%%%%%%%%%%%%%%%%%%%%%%%%%%%%%%%%%%%%%%

\chapter{Results and Discussion}

The NDVI and EVI maps generated from Landsat-8 and Sentinel-2 imagery indicated spatial variation in vegetation density within the study area.

The evapotranspiration maps generated using the Penman-Monteith method showed variation in crop water demand.

Crop coefficient maps developed using FAO crop coefficients helped estimate crop water requirements for paddy and ragi crops.

The comparison between evapotranspiration and crop yield revealed that regions with higher vegetation indices generally exhibited better crop performance and yield.

%%%%%%%%%%%%%%%%%%%%%%%%%%%%%%%%%%%%%%%%%%%%%%%%%%%%%%%%%%%%%%%%
%% Sample Table
%%%%%%%%%%%%%%%%%%%%%%%%%%%%%%%%%%%%%%%%%%%%%%%%%%%%%%%%%%%%%%%%

\begin{table}[H]

\centering

\caption{Sample Crop Yield Comparison}

\begin{tabular}{|c|c|c|}

\hline

Region & ET Value & Crop Yield \\

\hline

Region 1 & XX & XX \\

Region 2 & XX & XX \\

Region 3 & XX & XX \\

\hline

\end{tabular}

\end{table}

%%%%%%%%%%%%%%%%%%%%%%%%%%%%%%%%%%%%%%%%%%%%%%%%%%%%%%%%%%%%%%%%
%% Chapter 6
%%%%%%%%%%%%%%%%%%%%%%%%%%%%%%%%%%%%%%%%%%%%%%%%%%%%%%%%%%%%%%%%

\chapter{Conclusion}

The study successfully demonstrated the application of Remote Sensing and GIS techniques for estimating evapotranspiration and crop water requirements in the Hemavati Left Bank Canal command area.

NDVI, EVI, ET, and crop coefficient maps generated using Landsat-8 and Sentinel-2 imagery provided useful information regarding crop health and irrigation demand.

The study highlights the usefulness of QGIS and Remote Sensing techniques for sustainable irrigation planning and agricultural water management.

%%%%%%%%%%%%%%%%%%%%%%%%%%%%%%%%%%%%%%%%%%%%%%%%%%%%%%%%%%%%%%%%
%% Future Scope
%%%%%%%%%%%%%%%%%%%%%%%%%%%%%%%%%%%%%%%%%%%%%%%%%%%%%%%%%%%%%%%%

\section{Future Scope}

\begin{itemize}

\item Integration of machine learning techniques for crop yield prediction.

\item Use of high-resolution satellite imagery.

\item Real-time irrigation monitoring systems.

\item Multi-season crop analysis.

\end{itemize}

%%%%%%%%%%%%%%%%%%%%%%%%%%%%%%%%%%%%%%%%%%%%%%%%%%%%%%%%%%%%%%%%
%% Appendix
%%%%%%%%%%%%%%%%%%%%%%%%%%%%%%%%%%%%%%%%%%%%%%%%%%%%%%%%%%%%%%%%

\appendix

\chapter{Appendix}

Additional maps, tables, and calculations can be included here.

%%%%%%%%%%%%%%%%%%%%%%%%%%%%%%%%%%%%%%%%%%%%%%%%%%%%%%%%%%%%%%%%
%% Bibliography
%%%%%%%%%%%%%%%%%%%%%%%%%%%%%%%%%%%%%%%%%%%%%%%%%%%%%%%%%%%%%%%%

\backmatter

\begin{thebibliography}{99}

\bibitem{1}

Allen R G, Pereira L S, Raes D and Smith M,
``Crop Evapotranspiration Guidelines for Computing Crop Water Requirements,''
FAO Irrigation and Drainage Paper 56, 1998.

\bibitem{2}

Rouse J W, Haas R H, Schell J A and Deering D W,
``Monitoring Vegetation Systems in the Great Plains with ERTS,''
NASA Special Publication, 1974.

\bibitem{3}

QGIS Development Team,
``QGIS Geographic Information System,''
Open Source Geospatial Foundation, 2024.

\bibitem{4}

Jensen J R,
``Remote Sensing of the Environment,''
Pearson Education, 2007.

\end{thebibliography}

\end{document}
